\SourceLeaf{353}%
theoretical part. This has the further advantage that some investigations of interest in their own right, which are closely related to the others, prepare the way for them and shed more light on them, can be developed in greater detail than would have been consistent with treating the topics uniformly under the earlier plan. This holds even where those investigations had no direct application in the measurements in question.

The investigation that forms the subject of the present \textit{first} memoir belongs in particular to this class. The main part of the memoir consists of a method by which a triangulation network on the surface of an ellipsoid of revolution can be calculated as if it lay on a sphere, without any loss of accuracy. The method rests on the solution of a much broader problem. The author gave the solution in a memoir written in 1822 and brought into print by Conference Councillor \textsc{Schumacher} in the third fascicle of the \textit{Astronomische Abhandlungen}, under the title \textit{Allgemeine Auflösung der Aufgabe, die Theile einer gegebenen Fläche auf einer andern gegebenen Fläche so abzubilden, dass die Abbildung dem Abgebildeten in den kleinsten Theilen ähnlich wird.} For brevity, and because these representations of one surface on another are such a rich resource for calculations in higher geodesy that they well deserve a name of their own, the author has called those that meet the stated condition \textit{conformal mappings}. Thus the otherwise vague adjective is always to be understood here in a precisely defined sense. \textsc{Mercator}'s projection and the stereographic projection are familiar examples of conformal mappings of the sphere onto the plane.

It scarcely needs to be pointed out that similarity in the smallest (infinitesimal) parts must be carefully distinguished from similarity in all finite parts. The latter can be attained only in special cases, when the first surface can be developed onto the second itself or onto a surface similar to it. In general, however, when conformality consists \textit{only} in the similarity of the smallest parts, the \textit{scale ratio}, that is, the ratio of corresponding infinitesimal lines on the two surfaces to one another, is a number that %
